%% lampiran/B-hyperparameter-baseline.tex
%% Lampiran B — Hyperparameter Lengkap untuk Reproducibility

\chapter{Hyperparameter Lengkap dan Konfigurasi Pelatihan}
\label{lamp:hparam}

Lampiran ini memuat hyperparameter arsitektur dan konfigurasi pelatihan
lengkap untuk seluruh model yang dievaluasi pada Bab~\ref{bab:hasil-pembahasan},
demi memastikan reproducibility hasil eksperimen. Tabel~\ref{tab:lampB_tier_s}
merangkum tiga arsitektur Tier-S yang dipilih untuk evaluasi multi-\emph{seed}
penuh. Bagian~\ref{lamp:hparam_lstm}--\ref{lamp:hparam_informer} mendokumentasikan
arsitektur yang dikaji pada studi pendahuluan namun tidak diaktifkan sebagai
\emph{baseline} utama karena termasih dalam Tier-B atau Tier-C berdasarkan hasil
seleksi awal.

\begin{table}[ht]
  \centering
  \caption{Ringkasan tiga arsitektur Tier-S yang dievaluasi.}
  \label{tab:lampB_tier_s}
  \footnotesize
  \begin{tabularx}{\textwidth}{@{}l X r r@{}}
    \toprule
    \textbf{Model}      & \textbf{Paradigma Inti}                             & \textbf{Param.\ (PHM)} & \textbf{Param.\ (XJTU)} \\
    \midrule
    Mamba-xLSTM-Net     & SSM selektif + mLSTM + \emph{gated fusion}          & 898.481                & 811.713                 \\
    N-BEATS-xLSTM-RUL   & Basis polinomial + xLSTM \emph{temporal front}      & 459.592                & 457.672                 \\
    SparseGate-TCN-RUL  & TCN dilated + \emph{sparse gating}                  & 249.192                & 249.192                 \\
    \bottomrule
  \end{tabularx}
\end{table}

% -----------------------------------------------------------------
\section{Mamba-xLSTM-Net (Arsitektur Usulan)}
\label{lamp:hparam_proposed}
% -----------------------------------------------------------------

Mamba-xLSTM-Net merupakan arsitektur fokus penelitian ini. Hyperparameter
berikut berlaku untuk konfigurasi eksperimen utama (Bab~\ref{bab:hasil-pembahasan}).

\begin{table}[ht]
  \centering
  \caption{Hyperparameter arsitektur Mamba-xLSTM-Net.}
  \label{tab:lampB_mamba_xlstm_arch}
  \small
  \begin{tabular}{@{}lll@{}}
    \toprule
    \textbf{Hyperparameter} & \textbf{Nilai} & \textbf{Keterangan} \\
    \midrule
    \texttt{d\_model} & 128 & Dimensi representasi model \\
    \texttt{mamba\_d\_state} & 128 & Dimensi SSM \emph{state} Mamba \\
    \texttt{mamba\_expand} & 2 & \emph{Inner dim} = $2 \times$ \texttt{d\_model} \\
    \texttt{num\_mamba\_blocks} & 4 & Jumlah blok Mamba bidireksional \\
    \texttt{num\_mlstm\_heads} & 4 & Jumlah kepala mLSTM \\
    Input \emph{window} $L$ & 64 (PHM) / 32 (XJTU) & Panjang sekuens input \\
    Fitur input $F$ & 5 & Jumlah fitur \emph{band-energy} \\
    \texttt{dropout} & 0,1 & Dropout pada \emph{regression head} \\
    \texttt{head\_hidden} & 64 & Dimensi tersembunyi \emph{regression head} \\
    \bottomrule
  \end{tabular}
\end{table}

Arsitektur penuh terdiri dari: (1) proyeksi input $F \to$ \texttt{d\_model},
(2) $N_{\rm mamba}$ blok BiMamba (\texttt{d\_state = 128}, \texttt{expand = 2}),
(3) $N_{\rm xlstm}$ blok xLSTM (2 mLSTM + 1 sLSTM, konfigurasi \texttt{slstm\_positions = [1]}),
(4) \emph{gated fusion} antara output BiMamba dan xLSTM, dan
(5) \emph{regression head} (LayerNorm $\to$ Linear $\to$ GELU $\to$ Dropout $\to$ Linear $\to$ Sigmoid).
Rincian mekanisme \emph{gated fusion} dan \emph{bidirectional scan} disajikan
pada Bab~\ref{bab:dasar-teori} dan Bab~\ref{bab:metodologi}.

% -----------------------------------------------------------------
\section{N-BEATS-xLSTM-RUL}
\label{lamp:hparam_nbeats_xlstm}
% -----------------------------------------------------------------

N-BEATS-xLSTM-RUL merupakan arsitektur berbasis dekomposisi basis fisik
(\emph{trend}, \emph{wear}, \emph{shock}) yang digabungkan dengan modul
xLSTM sebagai \emph{temporal front-end}.

\begin{table}[ht]
  \centering
  \caption{Hyperparameter arsitektur N-BEATS-xLSTM-RUL.}
  \label{tab:lampB_nbeats_xlstm_arch}
  \small
  \begin{tabular}{@{}lll@{}}
    \toprule
    \textbf{Hyperparameter} & \textbf{Nilai} & \textbf{Keterangan} \\
    \midrule
    \texttt{hidden\_dim} & 96 & Dimensi tersembunyi setiap \emph{block} \\
    \texttt{trend\_blocks} & 2 & Jumlah blok basis polinomial tren \\
    \texttt{wear\_blocks} & 2 & Jumlah blok basis keausan (eksponensial) \\
    \texttt{shock\_blocks} & 2 & Jumlah blok impulsif \\
    \texttt{poly\_degree} & 4 & Derajat basis polinomial (\emph{trend blocks}) \\
    \texttt{model\_specific\_loss} & True & Aktifkan penalti monotonik per-komponen \\
    \bottomrule
  \end{tabular}
\end{table}

Setiap \emph{block} memproduksi estimasi komponen RUL yang terpisah
(\emph{trend}, \emph{wear}, atau \emph{shock}); hasil akhir adalah jumlah
ketiga komponen. Sifat dekomposisi ini memberikan interpretabilitas bawaan
(\emph{built-in interpretability}) per komponen degradasi tanpa memerlukan
analisis \emph{post-hoc}.

% -----------------------------------------------------------------
\section{SparseGate-TCN-RUL}
\label{lamp:hparam_sparsegate}
% -----------------------------------------------------------------

SparseGate-TCN-RUL adalah arsitektur TCN dilasi yang ditambahi mekanisme
\emph{sparse gating} untuk mendorong aktivasi \emph{channel} yang jarang
sehingga secara inheren lebih mudah diinterpretasi. Hyperparameter lengkapnya
ditampilkan pada Tabel~\ref{tab:lampB_sparsegate_arch}.

\begin{table}[ht]
  \centering
  \caption{Hyperparameter arsitektur SparseGate-TCN-RUL.}
  \label{tab:lampB_sparsegate_arch}
  \small
  \begin{tabular}{@{}lll@{}}
    \toprule
    \textbf{Hyperparameter} & \textbf{Nilai} & \textbf{Keterangan} \\
    \midrule
    \texttt{tcn\_channels} & [64, 64, 128, 128] & Lebar tiap lapisan TCN dilasi \\
    \texttt{tcn\_kernel} & 3 & Ukuran kernel konvolusi kausal \\
    \texttt{gate\_hidden} & 32 & Dimensi tersembunyi \emph{sparse gate} \\
    \texttt{attn\_d\_model} & 32 & Dimensi \emph{attention} ringan \\
    \texttt{attn\_heads} & 4 & Jumlah kepala \emph{multi-head attention} \\
    \texttt{dropout} & 0,1 & Dropout \\
    $\lambda_{\rm sparse}$ & 0,001 & Penalti \emph{L1} untuk sparsitas gate \\
    $\lambda_{\rm entropy}$ & 0,001 & Penalti entropi untuk diversitas gate \\
    \bottomrule
  \end{tabular}
\end{table}

% -----------------------------------------------------------------
\section{Konfigurasi Pelatihan Bersama (Semua Model)}
\label{lamp:hparam_training}
% -----------------------------------------------------------------

Semua model Tier-S dilatih dengan konfigurasi identik
(\texttt{configs/\allowbreak train/\allowbreak cloud\_full\_75.yaml})
untuk memastikan perbandingan yang adil. Perbedaan hanya pada
hyperparameter spesifik arsitektur yang didokumentasikan pada subbab
di atas.
Konfigurasi pelatihan lengkap yang berlaku untuk semua model ditampilkan
pada Tabel~\ref{tab:lampB_training}.

\begin{table}[ht]
  \centering
  \caption{Konfigurasi pelatihan bersama untuk semua arsitektur Tier-S.}
  \label{tab:lampB_training}
  \footnotesize
  \begin{tabularx}{\textwidth}{@{}l l X@{}}
    \toprule
    \textbf{Hyperparameter}              & \textbf{Nilai}                            & \textbf{Keterangan} \\
    \midrule
    \texttt{max\_epochs}                 & 75                                        & Jumlah \emph{epoch} pelatihan \\
    Optimizer                            & AdamW \citetitb{Loshchilov2019AdamW}      & \emph{Decoupled weight decay} \\
    \texttt{lr}                          & $8{,}0 \times 10^{-4}$                    & \emph{Learning rate} awal \\
    \texttt{weight\_decay}               & $3{,}0 \times 10^{-4}$                    & Koefisien \emph{weight decay} \\
    \texttt{xlstm\_lr\_mult}             & 0,75                                      & LR modul xLSTM = 0,75 $\times$ LR trunk \\
    \texttt{freeze\_xlstm\_epochs}       & 5                                         & xLSTM dibekukan 5 \emph{epoch} pertama \\
    \texttt{scheduler}                   & Cosine Annealing                          & $T_{\rm max} = 75$ \\
    \texttt{warmup\_epochs}              & 5                                         & Pemanasan \emph{cosine} di 5 \emph{epoch} awal \\
    \texttt{monotonicity\_weight}        & 0,05                                      & Bobot penalti monotonisitas RUL \\
    \texttt{gradient\_clip\_val}         & 1,0                                       & Klip gradien norma maksimum \\
    \texttt{early\_stopping\_patience}   & 9999                                      & Dinonaktifkan; 75 \emph{epoch} penuh \\
    \texttt{precision}                   & \texttt{bf16-mixed}                       & BF16 campuran (GPU Ampere+) \\
    Batch size (PHM2012)                 & 256 / 512                                 & 256 untuk Mamba + N-BEATS (paralel); 512 untuk SparseGate \\
    Batch size (XJTU-SY)                 & 256 / 512                                 & 256 untuk Mamba + N-BEATS; 512 untuk SparseGate \\
    \emph{Seed} acak                     & 42, 43, 44                                & Tiga \emph{seed} untuk estimasi stabilitas \\
    \bottomrule
  \end{tabularx}
\end{table}

Infrastruktur pelatihan: GPU NVIDIA A40 (46 GB VRAM), CUDA 12.8,
PyTorch 2.1 + PyTorch Lightning 2.1, framework manajemen konfigurasi Hydra 1.3,
pustaka \texttt{xlstm} 1.0 (NX-AI), \texttt{mamba-ssm} 2.2,
dan \texttt{causal-conv1d} 1.2.

% -----------------------------------------------------------------
\section{LSTM --- Studi Pendahuluan (Tidak Dilatih di Eksperimen Utama)}
\label{lamp:hparam_lstm}
% -----------------------------------------------------------------

Model LSTM berikut dikaji sebagai bagian dari studi pendahuluan seleksi
arsitektur (\emph{30-epoch screening}) namun tidak dilatih pada protokol
multi-\emph{seed} penuh karena berada di luar Tier-S berdasarkan hasil
seleksi. Hyperparameter ini dicantumkan untuk reproducibility studi
pendahuluan dan dapat dilihat pada Tabel~\ref{tab:lampB_lstm_arch}.

\begin{table}[ht]
  \centering
  \caption{Hyperparameter LSTM studi pendahuluan.}
  \label{tab:lampB_lstm_arch}
  \small
  \begin{tabular}{@{}lll@{}}
    \toprule
    \textbf{Hyperparameter} & \textbf{Nilai} & \textbf{Keterangan} \\
    \midrule
    \texttt{hidden\_dim} & 128 & Ukuran \emph{hidden state} LSTM \\
    \texttt{num\_layers} & 2 & Kedalaman lapisan LSTM \\
    \texttt{dropout} & 0,1 & Dropout antar lapisan \\
    \bottomrule
  \end{tabular}
\end{table}

% -----------------------------------------------------------------
\section{Transformer --- Studi Pendahuluan (Tidak Dilatih di Eksperimen Utama)}
\label{lamp:hparam_transformer}
% -----------------------------------------------------------------

Konfigurasi Transformer yang digunakan pada studi pendahuluan seleksi arsitektur
dirinci pada Tabel~\ref{tab:lampB_transformer_arch}. Arsitektur ini berada pada
Tier-B berdasarkan hasil \emph{screening} 30-\emph{epoch}.

\begin{table}[ht]
  \centering
  \caption{Hyperparameter Transformer studi pendahuluan.}
  \label{tab:lampB_transformer_arch}
  \small
  \begin{tabular}{@{}lll@{}}
    \toprule
    \textbf{Hyperparameter} & \textbf{Nilai} & \textbf{Keterangan} \\
    \midrule
    \texttt{d\_model} & 64 & Dimensi model \\
    \texttt{nhead} & 4 & Jumlah kepala \emph{attention} \\
    \texttt{num\_encoder\_layers} & 2 & Jumlah lapisan encoder \\
    \texttt{dim\_feedforward} & 256 & Dimensi FFN \\
    \texttt{dropout} & 0,1 & Dropout \\
    \bottomrule
  \end{tabular}
\end{table}

% -----------------------------------------------------------------
\section{Informer --- Studi Pendahuluan (Tidak Dilatih di Eksperimen Utama)}
\label{lamp:hparam_informer}
% -----------------------------------------------------------------

Konfigurasi Informer dengan \emph{ProbSparse attention} yang dikaji pada studi
pendahuluan dirinci pada Tabel~\ref{tab:lampB_informer_arch}. Informer
diklasifikasikan pada Tier-C berdasarkan \emph{screening} awal karena biaya
komputasi yang tinggi relatif terhadap kinerjanya pada dataset yang kecil.

\begin{table}[ht]
  \centering
  \caption{Hyperparameter Informer studi pendahuluan.}
  \label{tab:lampB_informer_arch}
  \small
  \begin{tabular}{@{}lll@{}}
    \toprule
    \textbf{Hyperparameter} & \textbf{Nilai} & \textbf{Keterangan} \\
    \midrule
    \texttt{d\_model} & 64 & Dimensi model \\
    \texttt{n\_heads} & 4 & Jumlah kepala \emph{ProbSparse attention} \\
    \texttt{e\_layers} & 2 & Jumlah lapisan encoder \\
    \texttt{d\_ff} & 256 & Dimensi FFN \\
    \texttt{dropout} & 0,1 & Dropout \\
    \texttt{factor} & 5 & Faktor \emph{ProbSparse sampling} \\
    \bottomrule
  \end{tabular}
\end{table}
